\subsection{Repairing auxiliary divisor support at one finite precision}

\begin{statement}{A.99}{Proved}
{A ray-class repair inside the actual rational norm fibre}
Retain the finite data $a,b,M,\ell,N,K,L,\mathcal B,\mathcal R$
of (A.95.1), with $\ell$ odd, and let
$\beta\in K^\times$ satisfy $\mathrm N_{K/\Q}\beta=b$.
Fix $r\ge1$ and a finite set $S$ of rational primes containing
two, the support of $b$, and every prime ramified in $L/\Q$.
Let $\mathfrak m$ be any finite ray modulus of $K$ supported
above $S\setminus\{2\}$.
Prescribe $\rho\in\operatorname{Gal}(\mathcal R/\Q)$ fixing
$\Q(\mu_M)$, as in A.95.
Then there is an actual element $\beta'\in K^\times$ and a
finite set $Q$ of rational primes, disjoint from $S$ and the
support of $\beta$, such that
\[
\begin{aligned}
 \mathrm N_{K/\Q}\beta'&=b,\\
 v_{\mathfrak p}(\beta')&=v_{\mathfrak p}(\beta)
                                  &&(\mathfrak p\mid S),\\
 \beta'/\beta&\equiv1\pmod{\mathfrak m},\\
 v_{\mathfrak p}(\beta')&\equiv0\pmod{\ell^r}
                                  &&(\mathfrak p\nmid S\cup Q).
\end{aligned}
 \tag{A.99.1}
\]
Every $q\in Q$ splits completely in both $L$ and
$\mathcal B$, and has cyclotomic Frobenius $\rho$ in
$\mathcal R$.
The set $Q$ may avoid any further prescribed finite set and
may lie above any prescribed bound.
At each such $q$ the valuation vector of $\beta'$ has
ordinary sum zero, so its rational norm contributes no
new coefficient prime.

This repairs one norm representative at the specified
finite precision.  For an intermediate field $K_0\subset K$,
the proof gives an explicit multiplier $u^{-t}$ and hence
the exact additional transition
\[
 \frac{\mathrm N_{K/K_0}\beta'}{\mathrm N_{K/K_0}\beta}
                         =(\mathrm N_{K/K_0}u)^{-t}.
 \tag{A.99.2}
\]
Its rational norm is one.  The statement does not assert
that (A.99.2) is an $\ell^r$th power in $K_0$, or that
separate applications preserve a previously chosen norm
chain in the radical tower.
\end{statement}
\begin{proof}
We use the primary ray class field with the exact Artin
map from ideals prime to $\mathfrak m$; see
\cite[V, Theorems 3.5--3.6]{MilneCFT}.
Write
\[
 C_{\mathfrak m}=\operatorname{Cl}_{\mathfrak m}(K)\{\ell\},
 \qquad H_{\mathfrak m}/K\text{ for its class field},
 \qquad F=L\prod_{\sigma\in G_\Q}\sigma(H_{\mathfrak m}).
 \tag{A.99.3}
\]
Then $F/\Q$ is Galois,
$V=\operatorname{Gal}(F/L)$ is an abelian $\ell$-group,
and $F/L$ is unramified outside the primes above $S$.
Since the ray modulus has no prime above two and $L/\Q$
is unramified at two, $F/\Q$ is unramified there.
Put $G=\operatorname{Gal}(F/\Q)$ and
$H=\operatorname{Gal}(F/K)$.
Restriction gives the actual surjection
$H\to C_{\mathfrak m}$ and the same exact quotients
$G/V=G_0$ and $H/V=U$ as in A.92.

For a rational prime $p\notin S$ define
$R_p^{\mathfrak m}$ as in (A.92.2), with the ray class
map to $C_{\mathfrak m}$.
All prime ideals involved are prime to $\mathfrak m$.
The proofs of A.92 and A.95 apply to these maps with
the following precise scope.
Every $p\notin S$ is unramified in $F$.
The split-prime character calculation, the extension
through $H^2(G_0,\Q_\ell/\Z_\ell)=0$, and the values
$f_i\chi(\operatorname{Frob}_p)$ therefore show that
$R_p^{\mathfrak m}$ is contained in the subgroup
generated by split-prime differences.
The field separation in A.95 uses only that
$F\mathcal R/\mathcal R$ is an $\ell$-extension
unramified at two, so it applies to (A.99.3) without
an unramifiedness assumption at the other controlled
primes.
Its same coset-character proof shows that the split
generators may all satisfy the prescribed $\rho$ and
split completely in $\mathcal B$.
In particular, one can choose a finite list of permitted
primes whose difference images generate any prescribed
finite list of these ray classes.
This is a statement about the displayed ray class
Artin maps; it does not replace the ray group by the
ordinary ideal class group.

Decompose the integral divisor of the given element as
\[
 \operatorname{div}(\beta)=D_S+D_{\rm aux},
 \qquad
 D_{\rm aux}=\sum_{p\notin S}
                   \sum_{\mathfrak p\mid p}v_{\mathfrak p}(\beta)
                                                   [\mathfrak p].
 \tag{A.99.4}
\]
These sums have finite support.
For each $p\notin S$, the local vector in $D_{\rm aux}$
has ideal norm zero, because
$v_p(\mathrm N\beta)=v_p(b)=0$.
Its class therefore lies in $R_p^{\mathfrak m}$.
The ray-class generation just proved supplies an
integral divisor $D_Q$ supported above finitely many
permitted primes $Q$, with norm zero separately at
each $q\in Q$, and satisfying
\[
                    [D_Q]=[D_{\rm aux}]
                                      \quad\text{in }C_{\mathfrak m}.
 \tag{A.99.5}
\]
The primes can avoid $S$, the finite support of $\beta$,
and any additional finite set.
The coefficients can be taken to be ordinary integers:
at a completely split prime the kernel is generated
by differences of prime basis vectors, and its image
is a finite group, so integer combinations give the
same image as $\Z_\ell$ combinations.

Let $D=D_{\rm aux}-D_Q$.
Its class in the full ray class group has order prime
to $\ell$ by (A.99.5).
There are consequently an integer $d>0$, prime to
$\ell$, and an actual element $u_0\in K^\times$ with
\[
              dD=\operatorname{div}(u_0),
                         \qquad u_0\equiv1\pmod{\mathfrak m}.
 \tag{A.99.6}
\]
The divisor $D$ is zero at every prime above $S$,
so $u_0$ is a unit at all these primes.
Its ideal norm to $\Q$ is zero, since both summands
of $D$ have norm zero.
Thus $\mathrm N_{K/\Q}u_0=\pm1$.
Replace $u_0$ by $u=u_0^2$ and $d$ by $2d$.
Because $\ell$ is odd, this new $d$ is still prime to
$\ell$, and we now have the exact identities
\[
       \operatorname{div}(u)=dD,\qquad
       \mathrm N_{K/\Q}u=1,\qquad u\equiv1\pmod{\mathfrak m}.
 \tag{A.99.7}
\]
This squaring keeps the ray condition; no multiplication
by a sign is used.

Choose an integer $t$ with $td\equiv1\pmod{\ell^r}$
and define $\beta'=\beta u^{-t}$.
Its rational norm is exactly $b$, and its local ratio
to $\beta$ is one modulo $\mathfrak m$.
Its divisor is the literal expression
\[
 \operatorname{div}(\beta')
           =D_S+(1-td)D_{\rm aux}+tdD_Q.
 \tag{A.99.8}
\]
The first term gives unchanged valuations above $S$.
The second term is divisible by $\ell^r$.
The last term is supported above $Q$ and has norm zero
at every one of its rational primes.
This proves every assertion of (A.99.1), including
the fixed rational norm and the support restriction.
Taking the actual intermediate field norm of the
formula $\beta'=\beta u^{-t}$ gives (A.99.2) and
establishes its stated scope.
\end{proof}

The ray modulus in this proof may prescribe any finite
local multiplicative precision away from two, while
the valuations above two are retained exactly.
Its exclusion of two is the explicit input that keeps
the auxiliary field unramified there and permits the
private-valuation separation from the $B$ radical.
The conclusion concerns valuations modulo $\ell^r$;
the divisor in (A.99.8) may still have nonzero
$\ell^r$-divisible coefficients at the old auxiliary
primes.  It is therefore a finite-precision repair,
with its effect on lower norm representatives given
exactly by (A.99.2).

\noindent The finite certificate
\path{code/norm_ray_support_certificate.py}, with retained output
\path{code/norm_ray_support_certificate_output.txt}, checks the
literal divisor coefficients, their weighted norms, the
prime-to-$\ell$ inverse exponents and the squaring operation
that preserves the ray congruence and makes the rational norm
exactly one.  The ray-class existence maps are proved above.
